\section{Exact enumeration: a rational function and a period-eight formula}
\label{sec:exact-odd-count}

The finite grammar can be made unambiguous by using the canonical
degree-two-pruned core.  This gives an exact count, strengthening the
quartic growth theorem.

\begin{theorem}[Exact unlabelled equality count]
There is one minimum-boundary two-tree at each of orders three,
four, and five, namely $K_3$, $K_4-e$, and the cone over $P_4$.
There are two at order six, seven at order seven, and exactly three
at every even order at least eight.  For every odd $n\ge9$,
\[
 a_n=\frac{n^4-20n^3+239n^2+c_rn+d_r}{192},\qquad r=n\pmod8,
\]
where
\[
\begin{array}{c|rrrr}
r&1&3&5&7\\ \hline
c_r&-1126&-1090&-1126&-1090\\
d_r&2250&1962&2298&1914
\end{array}
\]
In particular $a_n\sim n^4/192$ through odd $n$.
With $t=x^2$, the exact generating function is
\begin{equation}
 \sum_{\substack{n\ge7\\n\text{ odd}}}a_nx^n
 =\frac{x^7N(t)}{(1-t)^5(1+t)^2(1+t^2)},
 \label{eq:exact-count-gf}
\end{equation}
where
\[
 N(t)=7-3t-4t^2+8t^3-t^4+4t^5+11t^6-8t^7+t^8+t^9.
\]
\end{theorem}

\subsection{Canonical arms and their generating functions}

Use the state numbering of the compact odd grammar.  For
$j=0,\ldots,9$, let $F_j(x)$ count rooted component isomorphism
classes in state $j$, with $x$ marking vertices other than its two
root endpoints.  These are one-leaf components, so their pruned
core is a triangle path.  Their root endpoint triples are unequal
and fix an orientation.

There is a unique decomposition of such a rooted graph.  Retain
the two marked root endpoints and remove every other degree-two
vertex; this recovers its rooted core.  Its root triangle either
is the sole triangle, or has exactly one continuing core edge.
The ear multiplicities and the continuing edge determine the
last extension uniquely, including the predecessor orientation.
Iterate towards the leaf.  The only duplicate initial parameters
are $(p,0)$ and $(0,p)$ for $p=1,2,3$, which describe the same
canonically oriented cap and are counted once.  Thus the unary
transition equations are exact unlabelled equations, rather than
equations for derivations.  Their solutions are
\[
\begin{array}{c|l@{\qquad}c|l}
j&F_j(x)&j&F_j(x)\\ \hline
0&x^2&5&x^4(1+4t-3t^2)/(1-t)^3\\
1&2x^4/(1-t)&6&2x^5/(1-t)^2\\
2&x^3(1+4t-3t^2)/(1-t)^3&7&x^4\\
3&x^3&8&2x^4(1+t)/(1-t)\\
4&x^3(1+t)/(1-t)&9&x^4(1+t^2)/(1-t)^2
\end{array}
\]
These identities follow from all one-leaf rules with the three
initial seeds $F_0=x^2$, an added $x^3$ in $F_4$, and an added
$x^4$ in $F_9$.  Every extension increases order, so the solution
as formal power series is unique.

Define
\[
 A=F_0+F_1=\frac{x^2(1+t)}{1-t},\qquad
 B=F_3+F_4+F_6=\frac{2x^3}{(1-t)^2}.
\]
The $A$-arms have weight two and root degrees $(2,3)$; the
$B$-arms have weight three and root degrees $(2,4)$.  For an
unlabelled class $C$, unordered pairs have generating function
\[
 \operatorname{Pair}(C)(x)=\frac{C(x)^2+C(x^2)}2.
\]
Set
\[
 M=\operatorname{Pair}(A)
   =\frac{x^4(1+t+2t^2)}{(1-t)^2(1+t)}.
\]

\subsection{One or two branches}

If the core has one branch, its unique trivalent edge-node
canonically roots the graph.  Each arm has low root degree two
and high root degree at least three.  One root ear would make
the sum of final root degrees at least $3\cdot5-4+2=13$, so
there are no root ears.  No endpoint can be high in all three
arms, which would give degree at least seven.  Each therefore
has requirement at least six.  The degree sum is at least
eleven; degrees five and six would give one boundary and one
interior root, and weight at least six.  Hence both root
endpoints are interior of degree six.  The three arm
weights must be $2,2,3$.  The endpoint degree equations force
the weight-three arm to have root degrees $(2,4)$, and the
two weight-two arms to have their degree-three sides aligned
with its degree-two side.  Thus there are two $A$-arms and one
$B$-arm.  The unique $B$-arm fixes the root ordering, and the
$A$-arms form an unordered pair.  Every such combination has
requirement six at both final endpoints.  Therefore the exact
three-leaf contribution is
\begin{equation}
 U_3=x^2MB
 =\frac{2x^9(1+t+2t^2)}{(1-t)^4(1+t)}.
 \label{eq:count-u3}
\end{equation}

With two branches, the same root equations force two $A$-arms
at each branch and state 12 for the component leading to the
other branch.  The only entries into state 12 are an extension
of states 14 or 15, which are a pair of $A$-arms, and its
degree-six self-loop.  Hence the graph is determined by an
unordered pair of unordered pairs of $A$-arms, and a central
length $r\ge0$.  The core canonically recovers all these data.
Reversing the central strip exchanges its two branches: the
central core is the square of a path, with one ear on each
strictly internal consecutive edge, and the higher arm-degree
endpoint at the outer end of either terminal edge.  This
reversal preserves the endpoint roles for every $r$.
Consequently the exact four-leaf contribution is
\begin{equation}
 U_4=\frac{x^3}{1-x^2}\operatorname{Pair}(M)
 =\frac{x^{11}(1+4t^2+2t^3+6t^4+3t^6)}
 {(1-t)^5(1+t)^2(1+t^2)}.
 \label{eq:count-u4}
\end{equation}
The canonical decomposition proves that repeated or equal arms
are counted correctly; there is no assumption that grammar
derivations are unambiguous.

\subsection{Path cores and reflection}

For a path core, mark one of its two end triangles.  Its core
ear and opposite internal edge are intrinsic.  Exactly one
closed edge of this end triangle has ears, so this fixes the
root endpoint order.  The end component is a literal initial
cap in state 0,4, or9.  The terminal table with this cap first
and its orientation zero gives the end-rooted function
\begin{align*}
 E={}&x^4(F_2+2F_3+2F_4+F_6)\\
 &+x^5(3F_0+3F_1+F_5+F_7+F_8+F_9)\\
 &+x^6(2F_3+2F_4+F_6)\\
 ={}&\frac{4x^7(2+t-3t^2+t^3)}{(1-t)^3}.
\end{align*}
Selecting the other end is an involution.  Its fixed classes
are exactly graphs admitting an end-exchanging automorphism.
If their generating function is $S$, Burnside's elementary
two-element formula gives $U_2=(E+S)/2$.

The center of the core incidence path is an edge-node or a
triangle-node.  Its symmetry cases are exhaustive as follows;
$q$ counts ears on the central edge, or on each of the two
internal edges at a central triangle.
\[
\begin{array}{c|c|l|c}
\text{center}&\text{half states}&\text{matching}&q\\ \hline
\text{edge}&0,1&\text{same endpoint order}&1\\
\text{edge}&3,4,6&\text{opposite endpoint order}&1\\
\text{triangle}&0,1&\text{shared corner high}&0\\
\text{triangle}&3,4&\text{shared corner low}&0\\
\text{triangle}&3,4,6&\text{shared corner low}&1
\end{array}
\]
Here ``low'' and ``high'' refer to the half's two root degrees.
At an edge center the weight equation is
$2w+q+s_u+s_v=5$, so parity forces $q=1$; the root degree and
requirement tests give the first two rows.  At a triangle
center, let $u$ be the shared corner and $v,w$ the exchanged
corners.  Initially allow $r$ ears on the third edge.  Then
$D_u=2d_u+2q$, $D_v=D_w=d_v+q+1+r$, and the weight is
$2w+2q+r+s_u+2s_v$.  The degree bound gives $q\le1$.
If $r>0,q=1$, the shared corner would have degree six and
requirement at least eight, impossible.  If $r>0,q=0$ and
the shared corner is high, the other two corners are boundary
and the weight is at least $4+r+1\ge6$.  If the shared corner
is low, it is boundary.  Half high degree three then gives
weight at least eight; high degree four allows only $r=1$
and gives weight at least six; larger degrees are impossible.
Thus $r=0$.  The degrees reduce to
$D_u=2d_u+2q$ and $D_v=D_w=d_v+q+1$, and the weight equation is
$2w+2q+s_u+2s_v=5$.  If $u$ is low, its requirement is $5+q$,
and that at $v$ is $\max(A_v,5+q)$.  If $u$ is high, the
degree-six bound forces high degree three and $q=0$.  These
conditions give exactly the remaining three rows.

Isomorphic rooted halves always give the required reflection;
use an isomorphism in one direction and its inverse in the
other.  Conversely every end exchange identifies the halves
in one listed way.  Unequal root summaries preclude another
root-swapping identification within a half.  The center type
and its ear multiplicities are intrinsic, so the cases are
disjoint.  Doubling the half replaces $x$ by $x^2$, giving
\begin{align*}
 S_{\rm edge}&=x^3(F_0+F_1+F_3+F_4+F_6)(x^2),\\
 S_{\rm triangle}&=x^3(F_0+F_1+F_3+F_4)(x^2)
                  +x^5(F_3+F_4+F_6)(x^2),\\
 S&=\frac{2x^7(1+t-t^3)}{(1-t)^2(1+t)}.
\end{align*}
Thus
\begin{equation}
 U_2=\frac{x^7(5+6t-5t^2-5t^3+3t^4)}{(1-t)^3(1+t)}.
 \label{eq:count-u2}
\end{equation}
The path-core count is five at order seven and, for odd $n\ge9$,
equals $(n^2-25)/4+\lfloor(n+1)/4\rfloor$.

\subsection{Exact algebra and independent checks}

The complete-core exceptions contribute $2x^7$.  The other
classes have respectively two, three, and four core leaves,
so are disjoint.  Adding $2x^7+U_2+U_3+U_4$ yields
\eqref{eq:exact-count-gf}.

To certify the quasipolynomial without extrapolation, put
$m=(n-7)/2=4k+r$.  Its four Newton coefficient rows, indexed
by $r=0,1,2,3$, are
\[
\begin{array}{c|rrrrr}
0&8&109&502&896&512\\
1&18&177&678&1024&512\\
2&36&280&886&1152&512\\
3&66&420&1126&1280&512
\end{array}
\]
The corresponding polynomial is
$P_r(k)=\sum_{j=0}^4c_{rj}\binom{k}{j}$; subtract one only
at $m=0$.  Since
$\sum_{k\ge0}\binom{k}{j}z^k=z^j/(1-z)^{j+1}$, these rows
give a rational function with denominator $(1-t^4)^5$.
Cross-multiplication with \eqref{eq:exact-count-gf} is an
exact integer polynomial identity.  Expanding the four
polynomials in $n$ proves the displayed residue formula for
every odd $n\ge9$, with $a_7=7$ as its sole initial correction.

The standard-library checker \texttt{exact\_odd\_count.py}
verifies every one-leaf recurrence, terminal sum, symmetry
formula, unordered-pair identity, and the exact polynomial
certificate.  A separate program of the earlier session generates the complete
grammar by increasing order and canonicalizes graphs using
recursive ordered-edge incidence codes.  Equality of these
codes is exact equality of recursive tuples; rerooting at
all edges and both orders gives unlabelled isomorphism.
It reproduces every individual nauty class through order 13
and all stratified counts through order 35, where there are
$4686$ extremizers.  These experiments check the all-order
canonical proof rather than replace it.  The counting
correspondence is not represented as an end-to-end Lean
formalization.  The exact integer-polynomial identities for
the summed numerator and Newton coefficient formula are
separately kernel-checked in
\texttt{OddCountingIdentities.lean}.
